% =============================================================================
%  Notas del Profesor — Sesión 15: Paso por Referencia
%  Programación Avanzada — Universidad Panamericana
%  Compilar: pdflatex sesion15_paso_por_referencia_profesor.tex
% =============================================================================
\documentclass[12pt, oneside]{article}
\usepackage[profesor]{progavanzada}

\begin{document}

\portada{Sesión 15}{Paso por Referencia}{2026}

\tableofcontents
\newpage

% =============================================================================
\section{Objetivos de la sesión}
% =============================================================================

Al finalizar la sesión el alumno será capaz de:

\begin{enumerate}
  \item Predecir si una función modificará o no el argumento original, dado su encabezado.
  \item Declarar parámetros por valor, por referencia (\cpp{int\&}) y por apuntador (\cpp{int*}).
  \item Elegir el mecanismo de paso adecuado según el tipo de dato y la necesidad de modificación.
  \item Usar \cpp{const\&} para evitar copias costosas sin permitir modificaciones.
  \item Implementar funciones que "devuelven" múltiples valores mediante referencias.
  \item Implementar \cpp{swap} y explicar por qué necesita paso por referencia.
  \item Completar la implementación de Bubble Sort con las tres funciones separadas.
\end{enumerate}

\begin{notaprofesor}[Duración estimada]
  90 minutos (sesión simple). Distribución sugerida:\\
  15 min — apertura: la trampa invisible, demo en vivo \\
  20 min — paso por valor vs referencia, diagramas de memoria \\
  15 min — paso por apuntador, tabla comparativa, \cpp{const\&} \\
  15 min — múltiples resultados + arreglos \\
  25 min — Bubble Sort: \cpp{swap}, \cpp{comparar}, \cpp{bubblesort} en vivo \\
  \textit{Los ejercicios quedan como práctica para el cuaderno Jupyter.}
\end{notaprofesor}

% =============================================================================
\section{Prerrequisitos}
% =============================================================================

\begin{itemize}
  \item Apuntadores, referencias y aritmética de apuntadores (sesión 14)
  \item \textit{Array decay} y la pérdida del tamaño al pasar arreglos (sesión 14)
  \item Funciones: declaración, parámetros y valor de retorno (sesiones 11--13)
\end{itemize}

\begin{notaprofesor}
  Los alumnos ya vieron referencias en la sesión 14 como \textit{alias}. Esta sesión
  conecta ese concepto con el paso de parámetros, que es donde las referencias cobran
  utilidad práctica. Si algún alumno tiene experiencia en Python o Java, señala que en
  esos lenguajes los objetos se pasan por referencia automáticamente — el equivalente
  al \cpp{int\&} de C++ lo hacen transparente el intérprete o la JVM.
\end{notaprofesor}

% =============================================================================
\section{Secuencia de clase}
% =============================================================================

\subsection{Apertura — La trampa invisible (15 min)}

Escribe en el pizarrón (sin usar el cuaderno todavía):

\begin{codigo}
void sumaUno(int n) { n = n + 1; }

int a = 10;
sumaUno(a);
cout << a;
\end{codigo}

Pide a los alumnos que predigan el resultado \textit{por escrito} antes de
ejecutarlo. Luego ejecuta. El objetivo es que la disonancia entre lo esperado
(11) y lo obtenido (10) genere la pregunta; no expliques todavía.

\begin{notaprofesor}[Pregunta de arranque]
  ``¿Cuántos de ustedes esperaban 11?''\\
  Respuesta esperada: la mayoría levanta la mano.\\
  Respuesta correcta: 10. La función modifica su propia copia de \cpp{n}.\\
  Sigue con: ``¿Por qué C++ hace esto? ¿No sería más intuitivo que sí cambiara?''
  Deja que discutan 2 minutos y luego presenta el modelo de memoria en el pizarrón.
\end{notaprofesor}

Dibuja el diagrama de dos celdas separadas (la de \cpp{a} en \texttt{main} y la
de \cpp{n} en \texttt{sumaUno}) con direcciones distintas. Solo entonces introduce
la palabra ``copia'' y el término ``paso por valor''.

\subsection{Paso por referencia y por apuntador (35 min)}

Muestra las tres versiones del \cpp{sumaUno} en secuencia (valor, referencia,
apuntador) y pide que predigan el resultado de cada una antes de ejecutar.

\begin{notaprofesor}[Error muy frecuente: escribir \cpp{\&} en la llamada]
  Muchos alumnos escriben \cpp{sumaUno(\&a)} cuando la función ya es por
  referencia (\cpp{int\& n}). Esto es un error de compilación: \cpp{\&a} es
  de tipo \cpp{int*}, no \cpp{int\&}. El \cpp{\&} en la declaración del
  parámetro hace que el compilador tome la dirección \textit{implícitamente}.
  Enfatiza que el \cpp{\&} va \textbf{solo en la declaración}, no en la llamada.
\end{notaprofesor}

Presenta la tabla comparativa (valor / referencia / \cpp{const\&} / apuntador)
y el ejemplo de \cpp{const string\&}. Conecta con \cpp{cin >>} vs \cpp{scanf \&x}
para anclar el concepto en algo que ya conocen.

\begin{notaprofesor}[Pregunta de sondeo — const\&]
  ``Si \cpp{const\&} es eficiente y seguro, ¿por qué no usarlo para todo?''\\
  Respuesta esperada: para tipos primitivos la dirección (\cpp{int\&}) ocupa
  tanto como el dato mismo (4 vs 8 bytes en 64-bit), y la desreferencia añade
  una instrucción extra. El compilador puede optimizarlo, pero semánticamente
  es más claro \cpp{int} a secas cuando no se quiere modificar.
\end{notaprofesor}

\subsection{Múltiples resultados y arreglos (15 min)}

Demuestra la función \cpp{minMax} y pregunta: ``¿Por qué \cpp{minVal} y
\cpp{maxVal} son \cpp{int\&} pero \cpp{arr} es \cpp{const int*}?''
Deja que los alumnos respondan antes de explicar.

Cierra esta sección recordando el \textit{array decay} de la sesión anterior:
muestra que dentro de la función \cpp{sizeof(arr)} devuelve 8 (apuntador),
no el tamaño real del arreglo. La relación directa con la sesión 14 refuerza
la coherencia del curso.

\subsection{Bubble Sort (25 min)}

Construye el algoritmo en vivo, función por función:

\begin{enumerate}
  \item Escribe \cpp{swap} y verifica que intercambia dos variables. Pregunta:
        ``¿Qué pasaría si quitara el \cpp{\&}?'' (respuesta: intercambia copias, no cambia nada).
  \item Escribe \cpp{comparar}. Pregunta: ``¿Qué valor devuelve \cpp{comparar(3,5)}?
        ¿Y \cpp{comparar(5,3)}?'' antes de ejecutar.
  \item Escribe \cpp{bubblesort} y traza manualmente la primera pasada sobre el pizarrón
        con el arreglo de ejemplo \cpp{\{33, 42, 88, 3, 50, 27, 77, 69\}}.
\end{enumerate}

\begin{notaprofesor}[Demo en vivo — Inversión del comparador]
  Al terminar el Bubble Sort ascendente, cambia el \cpp{>} por \cpp{<} en
  \cpp{comparar} sin tocar ninguna otra función. Ejecuta y muestra que el
  arreglo queda en orden descendente. El punto: encapsular el criterio de
  comparación en una función separada hace el código flexible. Conecta con
  \cpp{std::sort} y los comparadores lambda.
\end{notaprofesor}

% =============================================================================
\section{Errores conceptuales frecuentes}
% =============================================================================

\begin{notaprofesor}[FAQ y errores comunes]

\textbf{1. ``¿Cuál es la diferencia real entre referencia y apuntador si los dos modifican el original?''}\\
La referencia es un alias que no puede ser nula ni reasignada; el apuntador es
una dirección que puede ser \cpp{nullptr} y puede apuntar a distintos objetos.
La referencia es más segura; el apuntador, más flexible. En paso de parámetros,
preferir referencia salvo que se necesite \cpp{nullptr} o aritmética de apuntadores.

\vspace{0.6em}
\textbf{2. Confundir \cpp{int\&} (referencia) con \cpp{\&x} (dirección-de)}\\
Un ejercicio útil: dar el mismo símbolo en tres contextos y pedir que identifiquen
el significado: \cpp{int\& r = x;} (declaración de referencia), \cpp{cout << \&x}
(dirección de \cpp{x}), y \cpp{int* p = \&x;} (se asigna la dirección de \cpp{x} a \cpp{p}).

\vspace{0.6em}
\textbf{3. Pensar que los arreglos se copian al pasar a funciones}\\
Viene de Java/Python donde los arrays sí se pueden copiar con facilidad. Demostrar
modificando el arreglo dentro de una función y mostrando que el original cambió.

\vspace{0.6em}
\textbf{4. Escribir \cpp{swap} sin variable temporal}\\
Algunos alumnos intentan \cpp{a = a + b; b = a - b; a = a - b;} o versiones con XOR.
Funcionan para enteros pero fallan para flotantes o tipos no numéricos. La versión
con temporal es la correcta y la que usa \cpp{std::swap} internamente.

\vspace{0.6em}
\textbf{5. Olvidar \cpp{n - i - 1} en el límite interno del Bubble Sort}\\
El índice interno puede ser \cpp{n - 1} y el algoritmo sigue siendo correcto,
solo hace comparaciones redundantes al final (los últimos \cpp{i} elementos
ya están ordenados). No es un error grave, pero es una oportunidad para hablar
de optimización.

\end{notaprofesor}

% =============================================================================
\section{Soluciones a los ejercicios}
% =============================================================================

\subsection{Ejercicio 1 — Elevar al cubo}

\begin{codigo}
void cubo(int& n) {
    n = n * n * n;
}

int arr[4] = {2, 3, 4, 5};
for (int i = 0; i < 4; i++)
    cubo(arr[i]);
// arr = {8, 27, 64, 125}
\end{codigo}

\begin{notaprofesor}[Errores típicos en este ejercicio]
  Algunos alumnos intentan \cpp{n = n\textasciicircum 3} (XOR bit a bit, no potencia).
  Recordar que C++ no tiene operador de exponenciación; se usa \cpp{pow(n, 3)} de
  \texttt{<cmath>} o la multiplicación explícita.
\end{notaprofesor}

\subsection{Ejercicio 2 — Tres resultados en una pasada}

\begin{codigo}
void estadisticas(const int* arr, int n,
                  int& minVal, int& maxVal, int& suma) {
    minVal = maxVal = arr[0];
    suma = arr[0];
    for (int i = 1; i < n; i++) {
        if (arr[i] < minVal) minVal = arr[i];
        if (arr[i] > maxVal) maxVal = arr[i];
        suma += arr[i];
    }
}

int datos[7] = {12, 5, 8, 21, 3, 17, 9};
int mn, mx, s;
estadisticas(datos, 7, mn, mx, s);
cout << "min=" << mn << " max=" << mx << " suma=" << s << endl;
// min=3 max=21 suma=75
\end{codigo}

\subsection{Ejercicio 3 — Normalizar un arreglo}

\begin{codigo}
void normalizar(double* arr, int n) {
    double maxVal = arr[0];
    for (int i = 1; i < n; i++)
        if (arr[i] > maxVal) maxVal = arr[i];

    for (int i = 0; i < n; i++)
        arr[i] /= maxVal;
}

double v[5] = {20.0, 50.0, 10.0, 80.0, 35.0};
normalizar(v, 5);
// v = {0.25, 0.625, 0.125, 1.0, 0.4375}
\end{codigo}

\begin{notaprofesor}
  Punto de discusión: ¿qué pasa si \cpp{maxVal} es cero? División por cero.
  Para un arreglo de datos reales habría que protegerse contra ese caso.
  Buen gancho para hablar de validación en los límites del sistema.
\end{notaprofesor}

\subsection{Ejercicio 4 — Bubble Sort descendente}

Solo hay que cambiar una línea en \cpp{comparar}:

\begin{codigo}
// Ascendente:
bool comparar(int a, int b) { return a > b; }

// Descendente (unica linea que cambia):
bool comparar(int a, int b) { return a < b; }
\end{codigo}

La explicación: \cpp{bubblesort} intercambia \cpp{arr[j]} y \cpp{arr[j+1]}
cuando \cpp{comparar} devuelve \cpp{true}. Si la condición de intercambio es
``el izquierdo es menor'', los elementos pequeños migrarán hacia el final
(orden descendente).

\begin{notaprofesor}[Extensión interesante]
  Pregunta abierta para los alumnos más avanzados: ``¿Podrías hacer que
  \cpp{comparar} sea un \textit{parámetro} de \cpp{bubblesort}, para que
  el usuario decida el criterio de ordenamiento?''
  \begin{codigo}
  void bubblesort(int* arr, int n, bool(*cmp)(int,int)) {
      for (int i = 0; i < n-1; i++)
          for (int j = 0; j < n-i-1; j++)
              if (cmp(arr[j], arr[j+1]))
                  swap(arr[j], arr[j+1]);
  }

  bubblesort(arr, n, comparar);
  \end{codigo}
  Este es el germen del polimorfismo de comportamiento, que en C++ moderno
  se expresa con lambdas y \cpp{std::function}. No es necesario profundizar
  ahora, pero despierta curiosidad hacia los temas futuros.
\end{notaprofesor}

% =============================================================================
\section{Conexiones con temas posteriores}
% =============================================================================

\begin{itemize}
  \item \textbf{Sesiones 16--17} — Recursividad: las funciones recursivas
        se llaman a sí mismas con paso por valor o por referencia según si
        necesitan acumular resultados.
  \item \textbf{Sesiones 18--20} — Memoria dinámica: \cpp{new} devuelve un
        apuntador que se pasa entre funciones; entender paso por apuntador
        es requisito.
  \item \textbf{Sesiones 21--25} — Orientación a Objetos: los métodos reciben
        \cpp{this} (apuntador al objeto) implícitamente; los constructores de
        copia usan \cpp{const T\&}.
  \item \textbf{C++ moderno} — \cpp{std::sort} con lambda, \cpp{std::swap},
        \cpp{std::move} (referencias rvalue, C++11) son extensiones directas
        del paso por referencia.
\end{itemize}

% =============================================================================
\section{Material de la sesión}
% =============================================================================

\begin{itemize}
  \item \textbf{Cuaderno:}
        \texttt{sesiones/sesion15\_paso\_por\_referencia.ipynb}
  \item \textbf{Ejemplos:}
        \texttt{materialInicial/sesion15/referencias01--06.cpp}, \texttt{sort.cpp}
  \item \textbf{Notas alumno:}
        \texttt{notas\_alumno/pdf/sesion15\_paso\_por\_referencia.pdf}
\end{itemize}

% =============================================================================
\section{Rúbrica de ejercicios}
% =============================================================================

\begin{center}
\begin{tabular}{p{6.5cm}cc}
  \toprule
  \textbf{Ejercicio / Criterio} & \textbf{Puntos} & \textbf{Evidencia} \\
  \midrule
  \multicolumn{3}{l}{\textit{Ejercicio 1 — Elevar al cubo}} \\
  \quad Parámetro declarado como \cpp{int\&}       & 4 & firma de la función \\
  \quad Resultado correcto para los 4 valores      & 4 & salida esperada \\
  \quad No usa \cpp{pow}                            & 2 & multiplicación explícita \\
  \midrule
  \multicolumn{3}{l}{\textit{Ejercicio 2 — Tres resultados}} \\
  \quad Tres parámetros de salida como \cpp{int\&} & 3 & firma correcta \\
  \quad Una sola pasada (un solo \cpp{for})         & 4 & código con 1 ciclo \\
  \quad Resultados correctos                        & 3 & min=3, max=21, suma=75 \\
  \midrule
  \multicolumn{3}{l}{\textit{Ejercicio 3 — Normalizar}} \\
  \quad Arreglo como \cpp{double*}, modifica el original & 3 & sin \cpp{const} en \cpp{arr} \\
  \quad Encuentra el máximo correctamente          & 3 & primera pasada \\
  \quad Divide todos los elementos                 & 4 & segunda pasada \\
  \midrule
  \multicolumn{3}{l}{\textit{Ejercicio 4 — Bubble Sort descendente}} \\
  \quad Identifica la línea correcta a cambiar     & 5 & solo \cpp{comparar} modificado \\
  \quad Explicación escrita correcta               & 5 & razonamiento con criterio de intercambio \\
  \midrule
  \textbf{Total}                                   & \textbf{40} & \\
  \bottomrule
\end{tabular}
\end{center}

\end{document}
